% Chapter 2

\chapter{Federated Learning Framework and Methodology}
\label{ch:chapter2}

Federated Learning is a distributed learning framework where multiple clients (devices or organizations) collaboratively train a shared global model under the coordination of a central server \cite{mcmahan2017communication}.

\section{Basic Workflow}

The FL process typically follows these steps:

\begin{enumerate}
    \item A global model is initialized at the server.
    \item The model is sent to selected client devices.
    \item Each client trains the model locally using its private data.
    \item Clients send updated model parameters (not raw data) back to the server.
    \item The server aggregates these updates to form a new global model.
    \item The process repeats over multiple rounds.
\end{enumerate}

This approach ensures that raw data never leaves the client device, which is the primary privacy advantage of FL.

\section{Types of Federated Learning}

Federated learning can be categorized based on how data is distributed:

\begin{itemize}
    \item \textbf{Horizontal FL:} Same features, different users (e.g., multiple banks with similar attributes but different customers).
    \item \textbf{Vertical FL:} Same users, different features (e.g., bank + e-commerce platform sharing complementary information).
    \item \textbf{Federated Transfer Learning:} Neither features nor users overlap significantly.
\end{itemize}

Each type introduces different technical and privacy challenges, particularly in terms of data alignment and secure computation \cite{zhang2021survey}.

\section{Why Federated Learning is Important}

\begin{itemize}
    \item Enables collaboration without violating privacy laws
    \item Reduces risk of centralized data breaches
    \item Supports edge devices (e.g., smartphones, IoT systems)
    \item Improves personalization while maintaining global learning
\end{itemize}

Despite these advantages, FL is not inherently secure---it only changes \textit{where} data resides, not \textit{how} information can be inferred.
